% Figaro Paper J: Scheduled Air Service as Coordinated Resource Markets
% Industrial / Systems Engineering register.
% Audience: industrial-engineering and systems-engineering scholars,
% supply-chain-coordination researchers, scheduled-service operations
% practitioners.
% Companion to Paper K (TradeLens replacement, container-shipping worked
% example) in the same discipline.
% Preprint, May 2026

\documentclass[11pt,a4paper]{article}

\usepackage[utf8]{inputenc}
\usepackage[T1]{fontenc}
\usepackage{amsmath,amssymb,amsthm}
\usepackage{mathtools}
\usepackage{booktabs}
\usepackage{graphicx}
\usepackage{hyperref}
\usepackage{cleveref}
\usepackage{enumitem}
\usepackage{xcolor}
\usepackage[margin=1in]{geometry}
\usepackage{natbib}
\bibliographystyle{plainnat}

\newtheorem{theorem}{Theorem}[section]
\newtheorem{proposition}[theorem]{Proposition}
\newtheorem{corollary}[theorem]{Corollary}
\theoremstyle{definition}
\newtheorem{definition}[theorem]{Definition}
\newtheorem{example}[theorem]{Example}
\theoremstyle{remark}
\newtheorem{remark}[theorem]{Remark}

\newcommand{\figaro}{\textsc{Figaro}}
\newcommand{\pay}{P}
\newcommand{\cumval}{G}

\title{%
  Scheduled Air Service as Coordinated Resource Markets:\\
  An Industrial-Engineering Reading of the Bonded Architecture%
}

\author{
  Alessandro Daliana\thanks{Corresponding author. \texttt{adaliana@figaro.org}}
}

\date{May 2026}

\begin{document}

\maketitle

% ════════════════════════════════════════════════════════════════════
\begin{abstract}
Scheduled passenger air service is coordination across a set of
resource markets: crew-hours and certifications, aircraft-time,
fuel volume, gate slots, catering counts, maintenance hours,
ground handling. Today these markets are internalized by
airlines, with the firm shape arising as the conventional
architecture's response to the high transaction and enforcement
costs of coordinating across markets without a firm-shaped
container. The bonded settlement primitive of the companion
mechanism paper changes those costs: enforcement is bilateral
and self-enforcing through asymmetric bonding, and atomic
resolution makes a multi-edge process tree resolvable from a
single signature. Under the bonded primitive the resource
markets become directly accessible to the passenger through
bonded commerce.

Two architectural properties carry the contribution:
\emph{asymmetric bonding within a single process} (the passenger
as \texttt{rootBuyer} posts $2P_i$ at each commit to each
resource provider, and each resource provider posts $2G$ against
the cumulative process value at the moment they commit; the
kernel enforces a single \texttt{rootBuyer} across every order in
the process) and \emph{buyer-initiated atomic resolution} (the
passenger settles the entire process at destination as a single
action). Together
they produce two consequences for scheduled-service coordination.
First, the resource markets become directly accessible to the
passenger: the firm-shaped airline becomes structurally optional
because the cross-market coordination it internalizes is
coordinatable through bonded commerce. Second, when operational
disruption occurs the seller cohort faces direct economic
incentive --- under locked bonds across every edge --- to
determine and pay compensation directly to the buyer, before any
external mechanism (specialized contracts, dispute forums,
insurance, regulators) engages. The contribution is not a better
airline; it is a different way of organizing scheduled air
service in which the firm-shaped container is structurally
unnecessary because the resource markets it internalizes are
directly coordinatable through bonded commerce. Cascading-delay
reduction follows as a consequence of bond posture and atomic
resolution; it is a property the architecture produces, not the
property the paper is for.
\end{abstract}

\medskip
\noindent\textbf{Keywords:}
scheduled-service coordination, supply-chain coordination,
process modeling, resource markets, weakest-link coordination,
industrial engineering

% ════════════════════════════════════════════════════════════════════
\section{Introduction}
\label{sec:intro}

Scheduled passenger air service is coordination across a set of
resource markets. The journey a passenger purchases is the joint
output of a flight crew (with required certifications and
duty-time bands), an aircraft and its operator, a quantity of
jet fuel, slot allocations at origin and destination gates,
per-passenger catering, a share of the day's maintenance
allocation for the operating airframe, and ground handling for
boarding and baggage. Each is a distinct resource the journey
draws on; each is supplied by a distinct set of operational
parties; each has natural units of measure already in industry
use --- crew-hours, type ratings, fuel volume in gallons or
kilograms, gate-slot time bands, meal counts, maintenance hours,
baggage counts. Air-traffic-control services are a further input
the flight consumes, but they are public-authority outputs not
market-traded; we list them only to bound the resource-market
inventory. Each resource market currently operates under its own
structural constraints (some auction-based, some bilateral, some
near-monopolistic at congested airports); the inventory below
makes no claim that those structures are uniform.

Today these resources are coordinated through airlines.
The airline contracts with crew agencies, fuel distributors,
airport authorities, catering firms, maintenance providers, and
ground-handling sub-contractors; sells tickets to passengers; and
absorbs the residual coordination cost. The firm shape is the
conventional architecture's response to a coordination problem
whose transaction and enforcement costs make firm-internalization
cheaper than market access for end customers
\citep{williamson1985}. The operations-research literature on
airline scheduling and operations
\citep{barnhart_cohn_2004,ball_etal_2007,lan_clarke_barnhart_2006}
largely takes this shape as given; analyses focus on optimization
within the firm-internal portfolio rather than on whether the
firm-internal organization is itself the right architecture.

The bonded settlement primitive of the companion mechanism paper
changes the calculus. Enforcement of bilateral commitments is
self-enforcing through asymmetric bonding; multi-party
coordination is resolvable from a single signature through atomic
resolution; the transaction-cost component that made
firm-internalization cheaper than direct market access becomes a
fixed capital lockup at the bonded edge. Under that change, the
resource markets that today's airlines internalize become
directly accessible to the passenger through bonded commerce. The
journey is composed as a process tree of bonded commitments
across those markets; each resource provider bonds against the
passenger's process; the passenger settles the entire tree at
destination through one atomic \texttt{resolveProcess} call.

\paragraph{Paper organization.}
\Cref{sec:primitive} summarizes the settlement primitive.
\Cref{sec:assembly} develops the architectural properties applied
to scheduled-service coordination: per-edge asymmetric bonding,
buyer-initiated atomic resolution, the two consequences for the
assembly, the kernel / composability-stack distinction, and where
specialized mechanism contracts can be developed for specific
resource markets. \Cref{sec:disruption} treats disruption
resolution: the seller-cohort compensation mechanism that arises
from buyer-dominance under bond pressure (the primary mechanism),
illustrated through a worked bond-posture example and through the
crew-allocation case where today's airline-internal auction
produces a documented misalignment that the bonded edge
eliminates, with the fallback mechanisms (specialized contracts,
coordination tokens, off-chain dispute forums, insurance,
regulators) treated downstream. \Cref{sec:schemas} treats the
schema requirements: one new schema, with the rest of the
assembly satisfied by existing infrastructure.
\Cref{sec:comparison} treats what the architecture preserves
(safety regulation, crew-labor protections, insurance,
competition law), what it changes (the cost-flow path and the
firm-internalization), and the non-cost reasons firms might
persist under the architecture nonetheless.
\Cref{sec:scope} states the scope conditions.
\Cref{sec:conclusion} concludes.

% ════════════════════════════════════════════════════════════════════
\section{The Settlement Primitive}
\label{sec:primitive}

The bonded primitive is developed formally in the companion
mechanism paper. We summarize the features the present argument
uses.

\paragraph{Asymmetric bonding (Mechanism~1).}
For a transaction with payment $\pay > 0$ and cumulative upstream
value $\cumval \geq \pay$, the buyer locks $2\pay$ and the seller
locks $2\cumval$. Cooperation is the unique profile surviving
iterated elimination of weakly dominated strategies. The
mechanism scales to $N$-party process trees through progressive
collateralization: each seller bonds against the cumulative value
of all upstream commitments, producing a mesh of independently
secured bilateral edges, each of which carries its own
equilibrium.

\paragraph{Buyer dominance with atomic resolution (Mechanism~2).}
Only the root buyer can trigger resolution, and resolution
settles all active orders in the process simultaneously or not at
all. The all-or-nothing rule induces a weakest-link subgame among
sellers: each seller's payout depends on universal cooperation
across the tree. The architecture produces cooperation pressure
of magnitude $\pay_i + 2\cumval_i$ on every other seller from
each seller's perspective, without explicit communication,
governance, or managerial layer.

\paragraph{What this implies for scheduled-service coordination.}
The bonded primitive's two mechanisms are exactly the two missing
properties in the contemporary scheduled-service architecture.
Mechanism~1 supplies a bilateral commitment whose enforcement
does not depend on either party's home jurisdiction's contract-law
regime; this is what makes it possible for the passenger's
commitment to a counterparty and that counterparty's commitment
to its sub-suppliers to be denominated in the same coordinative
unit. Mechanism~2 supplies the atomic resolution that propagates
a single commit-or-don't decision through the entire process
tree; this is what makes operational damages flow proportionally
to the unit that caused them rather than being absorbed at
whichever contract has the loosest cap. Neither mechanism is
novel to the scheduled-service domain; both are general-purpose
settlement properties whose application here is the present
paper's contribution.

% ════════════════════════════════════════════════════════════════════
\section{The Bonded Architecture Applied to Scheduled-Service}
\label{sec:assembly}

The substantive contribution of the bonded primitive to
scheduled-service coordination operates at the edge of the bonded
process tree, not at any particular position within it. Two
architectural properties supply the contribution; both are
independent of how the operational arrangement is organizationally
drawn.

\paragraph{Asymmetric bonding within a single process.}
The kernel enforces a strong structural rule: every order in a
process has the same buyer (the \texttt{rootBuyer}). The
passenger commits directly to each resource provider; there is
no intermediating coordinator between the passenger and any
seller. The first commit creates the process under the
passenger's address; every subsequent commit in the same process
extends the cumulative-value accumulator $G$ by the new payment.
At each commit the passenger posts $2P_i$ for that specific
payment, and the seller posts $2G_{\text{at commit time}}$ ---
the full cumulative process value at the moment they commit.
Bond posture is asymmetric: the passenger's per-commit bond
scales with the immediate payment, while each seller's bond
scales with the cumulative value of every upstream commit in the
process. Sellers committing later post correspondingly larger
bonds; the last seller bonds against the entire ticket value.

\paragraph{Buyer-initiated atomic resolution.}
The passenger calls \texttt{resolveProcess}, and the entire
process settles simultaneously. Every active order in the process
is resolved together or none is. The passenger's signature is the
only signature required for settlement, and atomicity guarantees
that no subset of orders settles without the rest.

\paragraph{Two consequences for the scheduled-service assembly.}
The two properties together produce two consequences.

First, \emph{the resource markets named in \Cref{sec:intro}
become directly accessible to the passenger through bonded
commerce.} Each resource provider --- crew operator, aircraft
operator, fuel supplier, gate operator, caterer, maintenance
provider, ground-handling firm --- can be a seller bonded against
the passenger's process, with payout depending on the passenger's
atomic resolution. Each provider's interests align with the
passenger's at the bonded commit rather than through the
firm-mediated incentive structure conventional airlines impose.
The misalignments today's firm-shaped container produces ---
between the airline's cost-minimization objective, the resource
provider's preferences, and the passenger's outcome --- dissolve
at the bond. We treat one such misalignment in detail in
\Cref{sec:disruption}: the crew secondary market.

Second, \emph{damage flow tracks bond posture across the
assembly's commits rather than the loosest contractual liability
cap.} Each seller carries asymmetric bond exposure proportional
to the cumulative process value at their commit; the seller-side
total grows super-linearly with the number of commits while the
passenger's posture stays at $2 P_{\text{ticket}}$. The
atomic-resolution rule operates on this set of bonds. When
something goes wrong, every seller's bond is locked across every
commit --- which is the structural pressure behind the
disruption-resolution mechanism we develop in
\Cref{sec:disruption}.

\paragraph{What the kernel constrains; what the composability stack carries.}
The kernel reads only per-order facts: buyer (which must be the
\texttt{rootBuyer} in every order in the process), seller,
payment, cumulative-value snapshot at commit, and the
atomic-resolution rule covering the process. The kernel-level
shape of a single process is therefore tightly constrained: one
buyer, one or more sellers, all under one cumulative accumulator,
all resolving atomically. Deeper compositions across multiple
processes (a seller in process $A$ becoming the
\texttt{rootBuyer} of process $B$ to coordinate their own
sub-suppliers) are admissible at the kernel layer but resolve
process-by-process, not across.

Above the kernel, the protocol-and-runtime composability stack
attaches schemas to specific orders (the
\texttt{scheduled-departure-v1} clause binds to whichever order
the parties want to subject to schedule-tolerance enforcement;
\texttt{handoff-v1} attestations attach to operational-handoff
orders), mechanism contracts operate on specific order patterns,
assembly specifications declare which schemas attach where and
which view definitions render which surfaces, and UIs render the
process per the assembly's declarations. An assembly's structure,
once specified, is a fixed verifiable artifact bound to its full
composability stack. The bonded-architecture argument of this
paper applies to whichever resource markets the designer composes
under one process, with the kernel-enforced structural rule
holding throughout.

\paragraph{Specialized mechanism contracts per resource market.}
Each resource market admits its own market mechanism. The
existing \texttt{DutchAuction} contract supplies one such
mechanism for cases where descending-price coordination is
appropriate (fuel procurement is a candidate where the supplier
market admits competition). Bilateral commerce-v1 commits cover
the cases where the parties contract directly. Future mechanism
contracts can be developed under the protocol-extension
discipline of the companion runtime-composition paper for
resources whose market mechanics deserve their own primitive: a
slot-allocation mechanism for gate-time bidding, a crew-dispatch
mechanism that respects type-rating and duty-time constraints, a
fuel-pool mechanism for shared ramp-side procurement among
co-located operators. The mechanism-contract space for scheduled
air service is underdeveloped today because the firm-shaped
container has absorbed the design choice; under direct market
access it becomes a first-class design surface.

% ════════════════════════════════════════════════════════════════════
\section{Disruption Resolution Under Buyer Dominance}
\label{sec:disruption}

We turn to the operational question: when something goes wrong
--- a delay, a slip, a service failure --- how is the disruption
resolved under the bonded architecture? The substantive answer is
that Mechanism~2 of the kernel (buyer dominance with atomic
resolution) creates structural pressure for the seller cohort to
determine and pay compensation themselves, directly to the buyer,
before any external mechanism engages. This is the primary
damage-management mechanism in the architecture. We treat it
first, place a bond-posture worked example, develop the
crew-allocation case as a concrete illustration, note the
IROps-scale extension, locate the result in the companion
mechanism paper's weakest-link result, and then position the
fallback mechanisms in their proper register downstream.

\subsection{The primary mechanism: seller-cohort compensation under bond pressure}

\paragraph{Bond posture per commit.}
Every order in a process shares the same \texttt{rootBuyer}; the
passenger as that \texttt{rootBuyer} posts $2P_i$ at each commit
to each resource provider, and each resource provider posts
$2G_{\text{at commit time}}$ where $G$ is the cumulative
process value at the moment of that commit. Sellers committing
later in the sequence post correspondingly larger bonds.
Aggregated across an $N$-seller process, the bonded exposure of
the seller cohort scales super-linearly with the number of
resource markets the assembly draws on.

\paragraph{A bond-posture worked example.}
Consider a stylized domestic short-haul passenger journey
totalling $\$250$. There is no airline-shaped coordinator: the
kernel enforces $\texttt{buyer} = \texttt{rootBuyer}$ in every
order in a process, so the passenger commits directly to each
resource provider. The seven commits below are all part of one
process under one $\texttt{rootBuyer}$ (the passenger). The
kernel maintains a single process-wide cumulative-value
accumulator $G$ that grows monotonically as each commit lands
($G \leftarrow G + P_i$). At every commit the buyer posts
$2P_i$ and the seller posts $2G_{\text{at commit time}}$. The
asymmetry shows up across the sequence: the passenger's
buyer-side bond at each commit scales only with the immediate
payment to that seller, while each seller's bond scales with
the full cumulative value of every upstream commit in the
process. Sellers committing later post correspondingly larger
bonds; the last seller bonds against the entire ticket value.
This is progressive collateralization in operation.
\Cref{tab:bond-posture} traces the sequence.

\begin{table}[h]
\centering
\small
\begin{tabular}{lrrrr}
\toprule
Commit (in sequence) & $P_i$ & $G$ after commit & Passenger bond $2P_i$ & Seller bond $2G$ \\
\midrule
Root: passenger $\rightarrow$ aircraft operator & $\$65$ & $\$65$ & $\$130$ & $\$130$ \\
Sub: passenger $\rightarrow$ fuel supplier & $\$45$ & $\$110$ & $\$90$ & $\$220$ \\
Sub: passenger $\rightarrow$ crew operator & $\$40$ & $\$150$ & $\$80$ & $\$300$ \\
Sub: passenger $\rightarrow$ maintenance provider & $\$15$ & $\$165$ & $\$30$ & $\$330$ \\
Sub: passenger $\rightarrow$ catering & $\$5$ & $\$170$ & $\$10$ & $\$340$ \\
Sub: passenger $\rightarrow$ ground handling & $\$15$ & $\$185$ & $\$30$ & $\$370$ \\
Sub: passenger $\rightarrow$ origin airport authority & $\$20$ & $\$205$ & $\$40$ & $\$410$ \\
Sub: passenger $\rightarrow$ destination airport authority & $\$20$ & $\$225$ & $\$40$ & $\$450$ \\
Sub: passenger $\rightarrow$ TSA (security screening) & $\$6$ & $\$231$ & $\$12$ & $\$462$ \\
Sub: passenger $\rightarrow$ federal aviation system & $\$19$ & $\$250$ & $\$38$ & $\$500$ \\
\midrule
Total bond locked across the process & --- & --- & $\$500$ & $\$3{,}512$ \\
\bottomrule
\end{tabular}
\caption{Stylized per-commit bond posture for a domestic
short-haul passenger journey under one process. The passenger
commits directly to each seller; there is no intermediating
coordinator (the kernel enforces buyer $= \texttt{rootBuyer}$ in
every order). $G$ is the process-wide cumulative-value
accumulator, monotonically growing with each commit. The
passenger's bond at each commit is proportional to the immediate
payment; each seller's bond is proportional to the cumulative
process value at the moment that seller commits. The last seller
to commit therefore posts the largest seller bond,
$2G_{\text{final}} = 2 \times \$250 = \$500$. Per-seller payments
are illustrative; absolute magnitudes vary with route, aircraft
type, and operating context.}
\label{tab:bond-posture}
\end{table}

Two features of the table carry the structural argument.

\emph{The asymmetry between buyer and seller bonds.} The
passenger's aggregate locked bond is $\$500$, which equals
$2 \times P_{\text{ticket}}$ --- the standard buyer posture. The
seller cohort's aggregate locked bond is $\$2{,}630$, more than
five times the passenger's posture, because each seller bonds
against the cumulative $G$ that has grown across the upstream
commits. This is not an accident of the example; it is the
asymmetric-bonding architecture in operation. As the assembly
adds more resource markets, the seller-side total grows
super-linearly with the number of commits while the passenger's
side stays at $2 P_{\text{ticket}}$.

\emph{The magnitude of the seller cohort's locked bond.} The
seller cohort's aggregate locked bond ($\$2{,}630$) is more than
ten times the ticket payment. A non-resolving passenger holds
that full pool locked across every commit in the process. The
pressure on the cohort to compensate the passenger and secure
resolution is structural in proportion that exceeds the ticket by
an order of magnitude.

The order in which sub-orders commit reorders which specific
seller posts the largest bond, but does not change the aggregate
seller cohort total: it depends only on the set of payments, not
the sequence. Whichever seller commits last bonds against the
full $G_{\text{final}} = P_{\text{ticket}}$.

\paragraph{Withholding resolution as the buyer's leverage.}
The passenger holds unilateral authority to call
\texttt{resolveProcess}. As long as the passenger does not call,
every seller's bond in the tree remains locked. Each seller faces
forfeit of their full $2G_i$ on continued non-resolution. The
buyer's atomic-resolution authority is therefore not a remedy in
the conventional sense (it does not pay the buyer anything beyond
bond recovery); it is leverage. The buyer can hold up the entire
tree's bond pool by simply not signing.

\paragraph{Sellers determine and pay compensation themselves.}
When operational disruption occurs and the buyer is left worse
off than the agreement contemplated, the seller cohort faces
direct economic incentive to make the buyer whole enough that the
buyer chooses to resolve. The mechanism unfolds operationally as
follows.

Sellers approach the buyer with compensation offers. Compensation
takes whatever form the buyer accepts: cash transfer, future
service credit, accommodation or meal vouchers, refund of the
ticket payment, alternative routing, partial or full refund
combined with services. The currency of compensation is at the
parties' discretion --- the kernel does not constrain it.

Sellers negotiate among themselves about which seller bears what
fraction of the compensation. The unit whose performance
materially failed bears more. Operational parties around the
failure who could not have prevented it (the catering provider
who had no role in a fuel slip; the gate operator whose service
was unaffected) contribute less or not at all. The negotiation is
real but routine: each seller knows their own bond exposure,
their counterfactual payout under cooperative resolution, and
their willingness-to-pay against losing the bond entirely.

The buyer accepts the compensation that satisfies them and calls
\texttt{resolveProcess}. Bonds release across the tree; every
seller recovers $2G_i + P_i$ less their share of the
compensation; the buyer recovers their bond plus the compensation
already paid. The settlement is bilateral and direct.

This happens fast. The negotiation occurs at the operational
boundary (the airport gate, the destination airport's customer
service desk, in some cases through the runtime UI's
disruption-coordination surface), not in a court filing six
months later. It happens before insurance claims, before
dispute-resolution forums, before regulatory engagement. Most
disruption resolves through this path because that is where the
bond pressure points.

\paragraph{The crew-allocation case.}
The mechanism is concrete in the crew market, where today's
firm-mediated arrangement produces a documented misalignment.
Airlines run bidding systems through which crew select trips by
seniority and preference. The bid systems are largely artifacts
of pilot-union and flight-attendant collective bargaining: senior
crew preserve seniority-allocation order through bargained
contracts, and the airline's discretion in the primary auction is
bounded by collective-agreement protections. Crew frequently view
the resulting allocations as imperfectly aligned with their
preferences and run secondary markets to trade trips with each
other after the primary auction settles. The secondary market is
operationally inefficient (it duplicates allocation work) and
structurally awkward (it requires the airline to permit and
process post-allocation swaps). Its existence is evidence that
the primary auction does not align crew interests with the
operational outcome that matters to passengers (qualified,
well-rested, willing crew on the flight).

Under the bonded architecture, the crew-allocation can operate as
a direct market between the passenger and the crew, without an
airline-internal primary auction or a post-allocation secondary
market at the passenger-coordination layer. A crew member bonds
against the passenger's process at the crew-edge; the crew
member's payout depends on the passenger's atomic resolution at
destination; the crew member's interests align with the
passenger's directly through the bonded edge. The
passenger-coordination auction and the secondary market both
become structurally unnecessary --- the bonded edge is the
passenger-side market, and the allocation discipline is whichever
mechanism contract the assembly designer composes for the
crew-market. The collective-bargaining institutions over wage
rates, duty-time protections, and seniority among crew remain;
they are labor-market institutions the bonded primitive does not
touch.

Under disruption --- say a crew member cannot fly --- the
crew-edge bond is locked alongside every other edge. The crew
operator and crew member negotiate compensation directly with the
passenger (find a substitute crew member at speed, accommodate
the passenger on the next flight, refund the leg, or whatever the
passenger accepts). Resolution follows compensation. The
conventional airline's role as intermediary in this negotiation
disappears.

\paragraph{IROps-scale coordination.}
The single-flight example understates the operational complexity
of large-scale irregular operations (IROps). A hub disruption
produces simultaneous bonded processes across hundreds of
flights, with thousands of bonded edges across the resource-market
mesh. The compensation negotiation at IROps scale is
operationally heavier than the per-flight case: each affected
passenger holds a separate bonded process with its own seller
cohort, and many resource providers (crew operators, aircraft
operators, ground-handling firms) appear in many trees at once.
We do not claim the bonded architecture eliminates this
coordination cost. We claim it routes the cost in a structurally
different direction. Under the conventional regime, the airline
absorbs the IROps coordination internally; sub-suppliers see only
their per-incident penalties; passengers see liability-cap
absorption. Under the bonded architecture, each affected
passenger faces a (busy, but bilaterally coordinated) seller
cohort negotiating their compensation directly. Operational
tooling for the resource provider --- whose bond is locked across
many simultaneous trees and who needs to coordinate compensation
responses at scale --- is itself one of the assembly-design
questions we identify as warranting empirical work.

\paragraph{Connection to the companion mechanism paper.}
The mechanism is the operational form of an analytical result
developed in the companion mechanism paper. Under buyer dominance
with atomic resolution, each seller bears direct economic
interest of magnitude $P_i + 2G_i$ in every other seller's
performance, because each seller's payout depends on universal
cooperation across the tree (the mechanism paper formalizes this
as a weakest-link subgame proposition: cooperation is the unique
equilibrium surviving iterated elimination of weakly dominated
strategies, induced by atomic resolution operating on the bonded
mesh). The proposition's analytical content is the formal
characterization of the pressure; the operational content is
sellers paying compensation to buyers directly to secure
resolution. The same structural shape produces the joint-liability
dynamic in Grameen-style microfinance, in which group members
chip in to cover a struggling member's installment before the
lender forecloses, because each member's continued credit access
is exposed to every other member's repayment. The bonded
architecture's seller-cohort-compensates-buyer dynamic is the
same structural mechanism scaled with chain depth and operating
without the social-substrate prerequisites Grameen requires.

\paragraph{Why this is structural, not procedural.}
The mechanism is not a courtesy that sellers may or may not
extend. It is a direct consequence of the bond posture under
atomic resolution. A seller who refuses to participate in
compensation negotiation watches their full bond forfeit
alongside every other seller's. The compensation mechanism
emerges from the architecture because the architecture has made
non-cooperation strictly worse than reasonable compensation for
every seller in the tree.

\subsection{Fallback mechanisms}

The primary mechanism resolves most disruption directly. Fallback
mechanisms engage when the primary does not, or when the
disruption involves a class of issue the primary cannot address.
We list them in the order of operational engagement.

\emph{Specialized mechanism contracts (protocol layer).} Where
the assembly composes a specialized mechanism contract --- a
delay-compensation contract that disburses pre-agreed amounts
under specified conditions, a slot-reallocation mechanism, a
fuel-pool mechanism --- the contract handles its specific
coordination automatically. The specialized contract complements
the primary mechanism by making specific compensation patterns
automatic rather than negotiated.

\emph{Coordination tokens (runtime layer).} Where the assembly
composes around a coordination token, the substrate's operational
state registers as token value. Token holders bear operational
quality through value impact rather than through direct
compensation. Token-value reflection is a diffuse signal
operating in parallel with the per-edge compensation mechanism,
not a substitute for it.

\emph{Off-chain dispute resolution forums.} The off-chain
agreement named in the bonded commitment binds a dispute forum
(state court, international arbitration center, ODR platform,
Schelling-point juror system). The forum engages when the
seller-cohort compensation does not resolve the dispute ---
typically because the dispute is substantive in a way bond
pressure cannot address (duress, mistake, frustration,
illegality, public-policy concerns). The bonded record is the
evidentiary input the forum operates on, with the Class~A /
Class~B taxonomy the companion evidence-and-law paper develops:
Class~A kernel byproducts (commit and resolution events) supply
the indisputable settlement record; Class~B schema-validated
attestations (\texttt{handoff-v1} timestamps,
\texttt{delivery-lifecycle-v1} stage events) supply the
substantive performance evidence.

\emph{Insurance compositions.} A passenger or operational party
who wants insurance against risks the bond pool cannot absorb
(catastrophic events, multi-day weather disruption, force-majeure
events) composes a parallel bonded process with an insurer. The
insurer's bond underwrites the risk; insurance settlement is its
own bonded event. Insurance composes with the bonded primitive
without modifying the kernel.

\emph{Regulators.} Aviation safety regulators, consumer protection
authorities, antitrust enforcement, and labor protection regimes
operate alongside the bonded architecture and engage in the
matters their authority covers. The architecture does not
displace them.

The substantive contribution of the bonded architecture to damage
management is that the primary mechanism --- seller cohort
negotiating and paying compensation directly to the buyer, before
external mechanisms engage --- is structural, immediate, and
operates without any of the fallbacks being needed in most cases.
The conventional regime's cascading-delay pathology --- damages
absorbed by whichever contract has the loosest liability cap,
with limited unit-level operational signal
\citep{lan_clarke_barnhart_2006} --- reduces under the bonded
architecture as a consequence of bond posture and atomic
resolution.

% ════════════════════════════════════════════════════════════════════
\section{Schemas}
\label{sec:schemas}

The assembly draws almost entirely on existing schema
infrastructure. One new schema is required: a binding from a
ticket commitment to a specific scheduled departure / arrival
pair, so the resolution path can verify whether the seller
delivered against schedule.

\paragraph{Existing schemas.}
\texttt{figaro-commerce-v1} carries the currency, payment, and
itemized line-items of every commitment in the assembly --- the
passenger ticket and every sub-procurement commitment.
\texttt{figaro-delivery-lifecycle-v1} carries staged progression
as a sequence of attestations; the schema's five-slot structure
can be re-used by reinterpretation (boarding initiated, cabin
ready, pushback, takeoff, arrival, or whatever projection the
assembly designer adopts), without authoring a new schema.
\texttt{figaro-handoff-v1} carries the operational handoffs
between operational parties.
\texttt{figaro-jurisdiction-v1} carries the off-chain
forum-selection clause for international flights, where the
parties want to specify a particular dispute-resolution forum.
\texttt{figaro-geo-v1} carries the origin and destination
geohashes where geographic-anchoring is operationally relevant.
The GHG family of schemas is available for assemblies that wish
to compose carbon-emission attestations into the ticket
commitment.

\paragraph{The new schema: \texttt{scheduled-departure-v1}.}
The schema binds a commerce-v1 ticket commitment to a specific
scheduled departure / arrival pair, so the resolution path can
verify whether the seller delivered against schedule. The fields
of the Layer~A spec are:

\begin{itemize}
  \item \texttt{flightDesignator} (string, IATA format, e.g.\ ``NK1234'')
  \item \texttt{scheduledDeparture} (ISO~8601 datetime, with explicit timezone)
  \item \texttt{scheduledArrival} (ISO~8601 datetime, same form)
  \item \texttt{originAirport} (string, ICAO~4-character, e.g.\ ``KORD'')
  \item \texttt{destinationAirport} (string, ICAO~4-character)
  \item \texttt{aircraftType} (string, optional, IATA aircraft code; omitted
    if the parties reserve substitution rights at the runtime tier)
  \item \texttt{toleranceSeconds} (integer, the agreed slip below which
    no penalty applies; e.g.\ 900 for a 15-minute window)
\end{itemize}

The schema enforces one structural constraint at the validator
contract: \texttt{scheduledArrival} must be strictly after
\texttt{scheduledDeparture}. We deliberately do not impose a
maximum on \texttt{toleranceSeconds}. Scheduled-service operations
vary widely in operationally meaningful tolerance windows, and
hard-coding a maximum at the validator contract would foreclose
legitimate use cases under the first-write-wins binding discipline
that makes the schemaId permanent. The contracting parties choose
the tolerance their operation supports; if a domain-specific
maximum is later wanted, it can be enforced at a domain-specific
schema (e.g.\ \texttt{figaro-passenger-departure-v1} as a tighter
sister of the present schema) without mutating the present
binding. The schema is the seed of a \emph{scheduled-service}
family; later extensions could cover scheduled ground transport,
scheduled deliveries with tight windows, and other
scheduled-binding use cases without reauthoring the family.

\paragraph{Schemas the architecture deliberately does not include.}
Three classes of schema are conspicuously absent. First, there is
no force-majeure schema that allows a counterparty to avoid bond
loss in weather events. Force-majeure handling is a
\emph{composition} concern, not a kernel concern: a party who
wants to insure against weather-driven delay composes a parallel
insurance process; the kernel neither knows about weather nor
adjudicates whether a given event qualifies. The companion
mechanism paper's no-escape-hatches analysis applies here
directly. Second, there is no schema that governs which delays
count toward bond loss; the \texttt{scheduled-departure-v1}
tolerance is the kernel's only such criterion, and any further
refinement is for the off-chain forum that adjudicates disputes
the parties bring to it. Third, there is no schema that grants
any party a unilateral reschedule right; rescheduling is a new
commitment with a new schemaId or a new instance, not a mutation
of the existing binding. Append-only identity is the discipline.

% ════════════════════════════════════════════════════════════════════
\section{Comparison with the Conventional Apparatus}
\label{sec:comparison}

The architectural change should be located precisely against the
existing apparatus to clarify what does and does not change.

\paragraph{What changes.}
The architecture changes the cost-flow path. Under the
contract-of-carriage regime, passenger schedule loss is bounded
by per-incident penalty caps and the airline's internal
willingness to issue voucher compensation; the supplier's
liability for sub-service slip is bounded by per-incident
contractual penalties; and the residual loss is absorbed by
whichever party's contract has the loosest cap. Under the bonded
architecture, schedule loss is allocated proportionally to the
party whose bond is exposed at the relevant process layer:
sub-supplier bond for sub-service slip, root-counterparty bond
for arrival-tolerance violation, passenger bond returned in full
on satisfactory delivery. The cost flow is direct, proportional,
and tracks the operational layer at which performance failed.

\paragraph{What does not change.}
Safety regulation is unchanged: the FAA in the US, EASA in the
EU, and analogous national civil-aviation authorities continue to
certify aircraft, pilot licensing, maintenance procedures, and
operational standards. Crew-labor protections are unchanged:
pilot-union and flight-attendant collective bargaining and
flight-time / duty-time limits under FAR Part 117 in the US (with
analogous regimes elsewhere) apply to the crew-supplier
sub-process under the bonded architecture exactly as they apply
under the conventional regime. Insurance and reinsurance are
unchanged in their underlying function; the bonded architecture
changes the risk profile aviation insurance markets price (a
smaller residual after bond allocation; concentrated risk at
whichever edges carry the largest cumulative-value bond posture;
a new class of bond-default risk at every operational party in
the tree) but does not displace the insurance function. Antitrust
and competition law are unchanged: operational-party markets
remain subject to anti-collusion enforcement, slot-allocation
regimes at congested airports remain governed by the relevant
authorities, and code-share arrangements remain subject to the
applicable competition regimes.

\paragraph{Non-kernel roles firms might still play.}
The bonded architecture removes the transaction-cost rationale
for firm internalization of scheduled-service coordination. It
does not remove other roles firms might play around the bonded
process. Three deserve mention. First, FAA Part 121 (and
analogous regimes in other jurisdictions) imposes a
single-certificate-holder requirement on commercial scheduled
passenger operations: the operating airline holds the air
carrier certificate, and many regulatory and safety functions
attach to that single legal entity. Under the bonded architecture
the certificate-holder is one of the resource-market sellers
(typically the aircraft-operator role); the certificate function
is preserved, but the certificate-holder is no longer a
cross-market coordinator. Second, knowledge accumulation across
operational cycles --- learning about reliability of specific
routings, weather patterns, sub-supplier reliability --- has
firm-level returns to scale that direct market access does not
naturally aggregate; firms retaining this knowledge are likely to
out-perform unaffiliated participants on quality, even within the
star-shaped bonded architecture. Third, brand identity functions
as a quality signal in markets with imperfect ex ante
observability of service quality. Brands cannot, however, become
kernel-level coordinators in the bonded architecture --- the
kernel forbids any party other than the passenger from being the
buyer of resource-market commits. A brand under the bonded
architecture is necessarily off-chain: a discovery surface, a
quality-signal aggregator, a multi-process curator that
recommends specific resource-provider compositions to passengers,
or a service-level guarantor that composes an insurance process
in parallel with the passenger's bonded process. None of these
roles re-introduce the firm-shaped cross-market coordinator the
architecture is designed to make structurally optional.

\paragraph{What this is not.}
The architecture is not a better airline. It is a re-architecture
of the cost-flow path in scheduled-service coordination. Whether
any specific operational party operates well or poorly within the
architecture is an operational-management question the
architecture is silent on.

The architecture is also not a passenger-protection regime. It is
a coordination protocol that has the side effect of routing
schedule-loss damages to the passenger more proportionally than
the contract-of-carriage regime does. Passenger-protection
regimes (the EU's Regulation 261/2004, the US's Department of
Transportation refund and compensation rules, and analogous
regimes elsewhere) continue to operate in their own register and
apply to the regulated entity in the relevant jurisdiction. The
bonded architecture's passenger-side benefit is structural rather
than regulatory.

% ════════════════════════════════════════════════════════════════════
\section{Scope}
\label{sec:scope}

The argument bounds itself by several scope conditions.

\paragraph{Asset specificity.}
The architecture works best where sub-suppliers compete in markets
with multiple available alternatives. Where a sub-supplier is a
near-monopolist (as airport authorities often are at hub airports
for gate operations, or as ATC services are by regulatory
construction), the buyer-side leverage at sub-procurement is
constrained, and the bonded architecture devolves toward the
conventional single-supplier regime. The companion
institutional-economics paper's discussion of asset specificity
applies here: the architecture is most effective in the
low-specificity region of the resource portfolio and least
effective at the high-specificity end.

\paragraph{Bond-currency access.}
The architecture presumes that all parties can post collateral in
a denomination the kernel accepts. For some sub-suppliers (e.g.\
small catering contractors operating in jurisdictions with
limited crypto on-ramps) this may be the binding operational
constraint. The companion statelessness paper's bond-currency
discussion applies in attenuated form: the constraint is real but
solvable through standard treasury and on-ramp infrastructure for
parties with the resources of typical airline sub-suppliers.

\paragraph{Aggregate liquidity at the operational tier.}
A high-volume operational party operating thousands of flights
per day under the bonded architecture would face aggregate
bond-custody requirements that scale with daily flight volume.
This is an operational treasury question: the party maintains a
working bond pool sized to its operational fleet, with bonds
released and rebound as processes resolve. The cumulative custody
at any moment is bounded by the active-process count multiplied
by per-process bond requirements; the architecture is
liquidity-intensive in steady state but not unbounded. The
financial-engineering question of how to size and finance the
bond pool is downstream of the architecture proposal here.

\paragraph{Network effects and adoption.}
The architecture's competitive position depends on adoption by a
sufficient fraction of sub-suppliers and by passengers willing to
trade the contract-of-carriage liability cap for the proportional
bonded architecture. We do not model the adoption dynamics. The
architecture is composable with the conventional regime: an
operational party can serve some passengers under the bonded
architecture and others under the contract-of-carriage regime,
and a sub-supplier can serve some counterparties under bonded
contracts and others under conventional ones. The transition is
not all-or-nothing.

% ════════════════════════════════════════════════════════════════════
\section{Conclusion}
\label{sec:conclusion}

Scheduled passenger air service is a coordination problem across
resource markets that the conventional architecture internalizes
inside firm-shaped airlines because the transaction and
enforcement costs of direct cross-market coordination were
prohibitive. The bonded settlement primitive collapses those
costs, and under the primitive the resource markets become
directly accessible to the passenger through bonded commerce. The
substantive contribution is not a better airline; it is a
different way of organizing scheduled air service in which the
firm-shaped container is structurally unnecessary because the
resource markets it internalizes are directly coordinatable
through bonded commerce. Cascading-delay reduction follows as a
consequence of bond posture and atomic resolution; it is a
property the architecture produces, not the property the paper is
for.

The architecture is composable with existing safety regulation,
crew-labor protections, insurance markets, and competition
regimes. It changes the cost-flow path without displacing the
regulatory and contractual layers that surround it. Whether the
architecture diffuses into commercial aviation is an adoption
question we do not adjudicate. What we do claim is that the
architecture exists, that it is well-defined, and that its
operational properties differ from the conventional regime in
ways the operations-research literature on airline scheduling has
been gesturing at without being positioned to deliver. The
companion mechanism paper supplies the equilibrium analysis; the
companion implementation paper supplies the verified Solidity
implementation; the present paper supplies the assembly that
turns those into a working scheduled-service coordination
architecture.

A natural next step is empirical: a pilot implementation of the
architecture for a domestic route with a manageable resource-market
portfolio, with measurement of bond-loss frequency,
disruption-resolution latency, and operational-improvement
outcomes at each resource provider, compared to the conventional
regime. The pilot would also surface the operational tooling
questions inherent to a star-shaped process the kernel enforces:
how passengers discover and select resource providers across the
seven (or more) markets the journey draws on, what off-chain
discovery / curation surfaces are useful, how cohort-internal
compensation negotiations play out at IROps scale, and how
multi-process composition operates when a resource provider
itself bonds with sub-suppliers under their own
\texttt{rootBuyer} role. We do not undertake the pilot here; we
identify it as the experimental work the architecture invites.

\section*{Acknowledgments}
The worked example originated in a project repository of
agent-coordination scenarios; the present paper formalizes the
scenario into an industrial-engineering register.

% ════════════════════════════════════════════════════════════════════
\begin{thebibliography}{99}

\bibitem[Barnhart and Cohn(2004)]{barnhart_cohn_2004}
Barnhart, C. and Cohn, A.
\newblock Airline Schedule Planning: Accomplishments and Opportunities.
\newblock \emph{Manufacturing \& Service Operations Management},
  6(1):3--22, 2004.

\bibitem[Ball et al.(2007)]{ball_etal_2007}
Ball, M.~O., Barnhart, C., Nemhauser, G., and Odoni, A.
\newblock Air Transportation: Irregular Operations and Control.
\newblock In C.~Barnhart and G.~Laporte, editors, \emph{Handbooks in
  Operations Research and Management Science: Transportation},
  volume~14, pages 1--67. Elsevier, 2007.

\bibitem[Lan et al.(2006)]{lan_clarke_barnhart_2006}
Lan, S., Clarke, J.-P., and Barnhart, C.
\newblock Planning for Robust Airline Operations: Optimizing Aircraft
  Routings and Flight Departure Times to Minimize Passenger Disruptions.
\newblock \emph{Transportation Science}, 40(1):15--28, 2006.

\bibitem[Williamson(1985)]{williamson1985}
Williamson, O.~E.
\newblock \emph{The Economic Institutions of Capitalism}.
\newblock Free Press, New York, 1985.

\bibitem[European Parliament and Council(2004)]{eu261_2004}
European Parliament and Council.
\newblock \emph{Regulation (EC) No 261/2004 establishing common rules
  on compensation and assistance to passengers in the event of denied
  boarding and of cancellation or long delay of flights}.
\newblock February 11, 2004.

\bibitem[FAA(2014)]{far117}
US Federal Aviation Administration.
\newblock \emph{14 CFR Part 117: Flight and Duty Limitations and Rest
  Requirements: Flightcrew Members}.
\newblock 2014.

\end{thebibliography}

\end{document}
